\section{Discussion}\label{sec:discussion}

The central lesson of this paper is that generic benchmark training does not transfer to an operating chemical plant: what a public construction-benchmark model delivers on plant imagery, and what an industrially adapted model delivers on the same frames, are different quantities, and the difference is the contribution. This section examines that design decision, adapting the detector to the environment it must watch, and closes with the limits of the present evaluation.

\subsection{Why Domain Adaptation Matters}

Public PPE corpora are dominated by construction imagery: near-frontal views, daylight, helmets and high-visibility vests \cite{roboflowppe,ppeyolo2024}. A phosphate-processing site violates each assumption: airborne dust lowers contrast, fixed CCTV cameras watch from high oblique angles at distances curated datasets rarely contain, and gas masks and goggles matter as much as helmets. The composite dataset answers these mismatches: fusion broadens pose, scale, and class coverage beyond any single source \cite{ahmad2024sh17}, the Roboflow gas masks corpus \cite{roboflowgasmasks} covers the class most specific to chemical processing, and weighted fine-tuning on the self-annotated OCP corpus closes the remaining site gap. The generic-versus-adapted comparison (Section~\ref{sec:results-comparison}) is the experiment this dataset was built for, and it is what separates the work from another YOLO variant: the evidence concerns what operating-plant imagery demands of training data, not the merits of a particular architecture.

\subsection{Limitations}

\begin{itemize}
  \item \textbf{Small and distant objects:} gloves and goggles occupy few pixels in wide CCTV views, and the nano-scale detector loses accuracy on them well before helmets or vests are affected.
  \item \textbf{Dust and weather:} heavy dust, rain, and low light plausibly push accuracy below the reported figures; robustness under these conditions has not been measured.
  \item \textbf{Gas-mask scarcity:} the class is deliberately not oversampled (Section~\ref{sec:methods-dataset}), so its effective training incidence is lower than OCP's or an earlier oversampled configuration would give it; augmentation adds visual variety but not additional distinct instances, and the class remains under-represented relative to its importance in chemical processing.
  \item \textbf{no\_goggle has no data:} despite occupying a slot in the fused ten-class map (Section~\ref{sec:methods-dataset}) and in the OCP annotation schema specifically, no\_goggle received zero labeled instances from any of the four sources in this run. It is not a thin class but an empty one. Its 0.0\% row in Table~\ref{tab:perclass} is consequently uninformative about detector capability; closing this gap requires deliberately annotating negative-goggle examples in a future OCP collection pass, not a training or evaluation change.
  \item \textbf{SH17 subsampling:} the caps on SH17's contribution (Section~\ref{sec:methods-dataset}: 1,500 training images, 300 each for validation and test) are a uniform random subsample, not stratified by class; classes already scarce within SH17 (e.g., boots, helmet) are thinned in the same proportion as abundant ones. Because the caps now also apply to validation and test, reported metrics are computed on a fixed 300-image random subsample of SH17's held-out splits rather than the full source, which was not previously the case; this trades some evaluation coverage for a wall-clock training budget that fits available compute. A class-aware cap that preserves rare-class density, and a check of whether the smaller SH17 eval subsample shifts reported metrics, are candidate refinements.
  \item \textbf{Comparison confounds:} the generic baseline is trained on far less data than the fused corpus, so the comparison measures what deploying the public benchmark model actually yields on plant imagery, not the isolated effect of domain coverage; a size-matched control is future work.
  \item \textbf{Camera placement:} the system inherits the blind spots of the existing CCTV layout; better detection cannot compensate for uncovered areas.
  \item \textbf{OCP imagery release:} the self-annotated site corpus and fine-tuning stage are complete, but public release of OCP frames remains restricted overall; Figure~\ref{fig:ocp-compare} shows one cleared frame as a qualitative illustration, not a representative sample, and the held-out OCP evaluation set itself remains unpublished.
\end{itemize}
